% from main.tex
%%%%%%%%%%  scan idx 486  |  folio 473  %%%%%%%%%%
où $i$ est une
immersion régulière, et $g$ un morphisme lisse. Comme $g$ est lisse, c'est un morphisme parfait (III 4.1). D'après VII 1.9, $i$ est également un morphisme parfait. Comme le composé de deux morphismes parfaits est parfait (III 4.5), on en déduit la proposition.

\textbf{Proposition 1.8.} Soient $f:X\to Y$ un morphisme d'intersection complète, $x$ un point de $X$, et une factorisation :
\[
\begin{tikzcd}[column sep=large,row sep=large]
U \arrow[r,"i"] \arrow[dr,"f"'] & V \arrow[d,"g"]\\
& Y
\end{tikzcd}
\]
où $U$ est un voisinage ouvert de $x$, $i$ une immersion régulière, et $g$ un morphisme lisse. On note $J$ l'idéal de $\mathcal O_V$ définissant $U$. Alors l'entier
\[
d_x=\operatorname{rang}\bigl((\Omega^1_{V/Y})_{i(x)}\bigr)-\operatorname{rang}\bigl((J/J^2)_x\bigr)
\]
est indépendant de la factorisation choisie de $f$.

Supposons que la restriction de $f$ à $U$ se factorise par deux $Y$-schémas lisses $V$ et $V'$, et posons $V''=V\times_Y V'$. Alors $V''$ est lisse sur $Y$, et aussi sur $V$ et $V'$. De plus, les immersions $X\to V$ et $X\to V'$ définissent une immersion $X\to V''$, régulière d'après 1.2. On peut donc se ramener à comparer les entiers $d_x$ relatifs à deux factorisations de $f$ par des $Y$-schémas lisses $V$ et $V'$ tels qu'il existe un diagramme commutatif au-dessus de $Y$ :
\[
\begin{tikzcd}[column sep=large,row sep=large]
& V' \arrow[d,"p"]\\
X \arrow[ur,"j"] \arrow[r,"i"'] & V,
\end{tikzcd}
\]
